\documentclass[../Main.tex]{subfiles}

\begin{document}
\section{Derivation Using the Interaction Picture}
Consider a Hamiltonian $H_0 = V(t)$ where $H_0$ is time-independent with known eigenstates and eigenvalues. Let $V(t)$ be small (in some sense), and we aim to calculate the effect of $V(t)$ order-by-order, and to find the probabilities of transitions between eigenstates of $H_0$.

Starting in the \schr picture, we have the \schr equation:
\begin{equation*}
    i\hbar \frac{\partial}{\partial t} \tstate = H\tstate = (H_0 + V(t)) \tstate
\end{equation*}
and so define the interaction picture by:
\begin{align*}
    \istate &= e^{iH_0 t / \hbar} \tstate \\
    \overline{A(t)} &= e^{iH_0 t / \hbar} A e^{-iH_0 t / \hbar}
\end{align*}
which moves some of the time evolution from states to operators (but only the $t$ dependence of $H_0$).

The remaining $t$ dependence in states is now:
\begin{align*}
    i\hbar \frac{\partial}{\partial t} \istate &= e^{iH_0 t / \hbar} \left(-H_0 + i\hbar \frac{\partial }{\partial t}\right) \tstate \\
    &= e^{iH_0 t / \hbar} V(t) \tstate \\
    &= e^{iH_0 t / \hbar} V(t) e^{-iH_0 t / \hbar}e^{i H_0 t / \hbar} \istate \\
    &= \overline{V(t)}~\istate
\end{align*}
This can be rewritten as an integral equation:
\begin{equation*}
    \istate = \oket{\psi(0)} -\frac{i}{\hbar} \int_{0}^{t} \overline{V(s)}~\oket{\psi(s)} ds 
\end{equation*}
which is useful for generating a series approximation to $\istate$. For example, by repeatedly substituting in our approximations for $\istate$, starting with $\oket{\psi(0)}$, 
\begin{align*}
    \istate &= \oket{\psi(0)} - \frac{i}{\hbar} \int_{0}^{t} \overline{V(s)}~\oket{\psi(0)} ds \\
    &\qquad+ \left(-\frac i \hbar\right)^2 \int_{s = 0}^{t} \int_{u = 0}^{s} \overline{V(s)}~\overline{V(u)}~\oket{\psi(u)} du~ds
\end{align*}
which gives the correct expression to order $V$ and shows that the next term is of order $V^2$ (we cannot evaluate this last double integral because we do not know $\oket{\psi(u)}$).
\section{Transition Probabilities}
Let $\ket a$ and $\ket b$ be eigenstates of $H_0$, with eigenvalues $E_a$ and $E_b$ such that $\inn a b = 0$. A transition from the state $\ket{\psi(0)} = \oket{\psi(0)} = \ket a$ to a state $\ket b$ at time $t$ occurs with probability:
\begin{equation*}
    \abs{\inn{b}{\psi(t)}}^2 = \abs{\langle b \istate}^2
\end{equation*}
since we are working with picture which give the same results for observable quantities.

To order $V$, the amplitude we need is:
\begin{align*}
    \langle b \overline{\mid \psi(t)\rangle} &= \inn b a - \frac i \hbar \int_{0}^{t} \opbraket b {\overline{V(s)}} a ds \\
    &= -\frac i \hbar \int_{0}^{t} \opbraket{b}{e^{iH_0 s / \hbar} V(s) e^{-iH_0 s / \hbar}}{a} ds \\
    &= -\frac i \hbar \int_{0}^{t} e^{i(E_b - E_a) s / \hbar} \opbraket{b}{V(s)}{a} ds 
\end{align*}
and so the probability for a transition is:
\begin{equation*}
    \frac1{\hbar^2} \abs{\int_{0}^{t} e^{i\omega s}\opbraket{b}{V(s)}{a} ds}
\end{equation*}
correct to order $V^2$, where $\hbar \omega = E_b - E_a$.

An important special case is when $V(s)$ is constant (at least on $s \in [0, t]$). Then the amplitude is:
\begin{equation*}
    -\frac i \hbar \int_{0}^{t} e^{i\omega s} \opbraket b V a ds = \frac1{\hbar\omega} (1 - e^{i\omega t}) \opbraket b v a.
\end{equation*}
The probability of a transition, to order $V^2$, is:
\begin{equation*}
    P(t) = \frac1\hbar \left(\frac{\sin(\omega t / 2)}{\frac12 \omega}\right)^2 \abs{\opbraket b V a}^2
\end{equation*}
and so for fixed $\omega \neq 0$ we have a periodic function of $t$.

Note:
\begin{equation*}
    P(t) = \frac{t}{\hbar^2} f_t(\omega) \abs{\opbraket b V a}^2
\end{equation*}
where $f_t(\omega) = \frac1t \left(\frac{\sin(\omega t / 2)}{\omega / 2}\right)^2$. This is a sine-squared function modulated by a $\frac1t$ curve.

We see that:
\begin{itemize}
    \item $\lim_{\omega \to 0} f_t(\omega) = t$
    \item $\int_{-\infty}^{\infty} f_t(\omega) d\omega = 2\pi$
    \item As $t$ increases, $f_t(\omega)$ is concentrated more and more around $\omega = 0$. Therefore, $P(t)$ becomes non-negligible only when $\omega \approx 0$, so when $E_b \approx E_a$.
\end{itemize}
\section{Fermi's Golden Rule}
The total probability for a transition to a band of final states is:
\begin{equation*}
    P(t) = \sum_b \frac{t}{\hbar^2} f_t\left(\frac{E_b - E_a}{\hbar}\right) \abs{\opbraket b V a}^2
\end{equation*}
In applications we encounter very closely-spaced final states, or even a continuum.

Consider final states specified by a set of parameters but with $E_b$ variable, and define the density of states $\rho(E_b)$ such that $\rho(E_b) \Delta E$ to be the number of states with specified parameters and energy in range $[E_b, E_b + \Delta E]$. The probability of transition to some state in band $B$ is then given by replacing the sum over $b$ to the integral over the band.

If there are other unspecified parameters in the final state, then\\$\abs{\opbraket b V a}^2$ must depend on these only through $E_b$.

Our previous remarks about $f_t$ shows that we can approximate it, for large times, by $2\pi \delta(\omega)$. Substituting this into the integral form for $P(t)$ gives:
\begin{equation*}
    P(t) = \frac{t}{\hbar^2} \int_B 2\pi \hbar \delta(E_b - E_a) \abs{\opbraket b V a}^2 \rho(E_b) dE_b
\end{equation*}
The delta function means that $P(t)$ is zero if $E_a$ is not in the band. If it is in the band,
\begin{equation*}
    P(t) = \frac{2\pi t}{\hbar} \rho(E_a) \abs{\opbraket b V a}^2
\end{equation*}
\begin{remark}
    \label{remTimeBounds}
    This only works for $t$ ``not too small'' (in the sense that the delta function approximation works -- we require that changes in the function $\rho(E_b) \abs{\opbraket b V a}^2$ are small for changes in $E_b$ of order $\frac\hbar t$) and for t ``not too large'' (in the sense that this is still a small perturbation -- we require $P(t) \ll 1$).
\end{remark}

Under the conditions of remark~\ref{remTimeBounds}, we see that the transition probability depends linearly on time. That is, we have a constant transition rate:
\begin{equation*}
    \frac{P(t)}{t} = \frac{2\pi}{\hbar} \rho(E_a) \abs{\opbraket{b}{V}{a}}^2
\end{equation*}
which is \underline{Fermi's Golden Rule}.
\section{Interaction of Electromagnetic Fields with Atoms} %TODO: Fix overfull hbox
\subsection{Classical Electromagnetic Fields as a Perturbation}
We consider a hydrogen atom, with Hamiltonian $H_0$, and a time-dependent potential:
\begin{equation*}
    V(t) = e \opvec x \cdot \vec{\EField} \cos(\omega_0 t)
\end{equation*}
This represents a classical EM wave with frequency $\omega_0$, amplitude and direction given by $\vec\EField$ constant.

The first-order transition amplitude between atomic states $a$ and $b$ is then:
\begin{align*}
    &-\frac i \hbar \int_0^t e^{i(E_b - E_a)s / \hbar} \opbraket b {V(s)} a ds \\
    &= -\frac{i e}{\hbar} \int_0^t e^{i(E_b - E_a)s / \hbar} \cos(\omega_0 s) \opbraket{b}{\vec{\EField} \cdot \opvec{x}}{a} ds
\end{align*}
and we can analyse this in two parts: first the time dependence and then the matrix element $\opbraket{b}{\vec{\EField}\cdot\opvec{x}}{a}$.

For the time dependence, we separate $\cos(\omega_0 t)$ into two terms:
\begin{equation*}
    \int_0^t e^{i(E_b - E_a)s / \hbar} e^{\mp i\omega_0 s} ds
\end{equation*}
Then evaluating and taking the modulus squared of the result gives the contribution to probability:
\begin{equation*}
    \frac{\sin^2\left((\omega \mp \omega_0) t / 2\right)}{(\omega \mp \omega_0)^2 / 4}
\end{equation*}
together with a cross term that is small. Here $\omega = (E_b - E_a) / \hbar$. As $t$ becomes sufficiently large, these terms are sharply peaked around $\omega \approx \pm \omega_0$.

For $t$ sufficiently large, then, we have two kinds of transitions, $a \to b$ with $E_b = E_a \pm \hbar \omega_0$ which occur with equal probabilities.
\nonexaminablesubsection{Quantum Electromagnetic Fields, In Outline}
Quantisation of the electromagnetic field can be thought of in terms of the harmonic oscillator. For one mode of frequency $\omega_0$, the electric field has a Heisenberg picture operator of the form:
\begin{equation*}
    a e^{-i\omega_0 t} + a^\dagger e^{i\omega_0 t}
\end{equation*}
where $a$ and $a^\dagger$ destroy and create photons of frequency $\omega_0$.

Using this, in the above discussion we have:
\begin{align*}
    \ket a &= \ket{n~l~m} \ket N \\
    \ket b &= \ket{n'~l'~m'} \ket N'
\end{align*}
which are products of atomic states and photon states. As in the classical case, we have two terms with frequencies $e^{\mp i \omega_0 t}$ but these are now accompanied by factors:
\begin{align*}
    \opbraket{N'} a N &=
    \begin{cases}
        \sqrt{N} & N' = N - 1 \\
        0 & \text{otherwise}
    \end{cases} \\
    \opbraket{N'}{a^\dagger}N &=
    \begin{cases}
        \sqrt{N+1} & N' = N + 1 \\
        0 & \text{otherwise}
    \end{cases}
\end{align*}
Then upon taking the modulus squared we find rates for emission and absorption of photons are equal apart from factors $N$ and $N+1$. Here $N$ is the number of photons and relates to the intensity of the electromagnetic radiation. The $N$ factors relate to stimulated emission/absorption and the extra $1$ in the emission term relates to spontaneous emission which is a new effect due to quantisation.

For general atomic energies $E_i$ and $E_j$, these have a transition if $E_i - E_j = \hbar \omega_0$. The transition rates are of the form $A_{i\to j} + B_{i \to j}$ for emission rate and $B_{j \to i}$ for absorption rate.
\subsection{Selection Rules}
We investigate the time-independent matrix element,
\begin{equation*}
    \vec{\EField} \cdot \opbraket{b}{\opvec{x}}{a}
\end{equation*}
which must be non-zero for a transition to occur. The probability, or transition rate, is then proportional to $\abs{\vec\EField}^2$, which is the intensity or energy density in classical EM radiation.

Consider states of the hydrogen atom, $\ket a = \ket{n~l~m}, \ket b = \ket{n~l~m}$. Then the information on whether $\opbraket{b}{x_i}{a}$ comes from properties with respect to reflections and rotations. Recall that the parity operator $U$ satisfies $U^2 = 1$ and:
\begin{equation*}
    U \ket{n~l~m} = (-1)^l \ket{n~l~m},\qquad U\opvec{x} U = -\opvec x
\end{equation*}
Therefore,
\begin{align*}
    \opbraket{n'~l'~m'}{\hat{x}_i}{n~l~m} &= \opbraket{n'~l'~m'}{UU\hat{x}_iUU}{n~l~m} \\
    &= (-1)^{l + l' + 1} \opbraket{n'~l'~m'}{\hat{x}_i}{n~l~m}
\end{align*}
and so the matrix element is zero unless $l+l'$ is odd. This is the first selection rule.

To study the implications of angular momentum, we first show that $\hat x_i \ket{n~l~m}$ are like product states for addition of angular momentum $l$ and $1$.

Recall that $\cmmt{L_i}{\hat x_j} = i\hbar \epsilon_{ijk} \hat x_k$ defines a vector operator. We then define:
\begin{align*}
    X_1 &= -\frac1{\sqrt{2}} \left(\hat x_1 + i \hat x_2\right) \\
    X_0 &= \hat x_3 \\
    X_{-1} &= -\frac1{\sqrt{2}} \left(\hat x_1 - i \hat x_2\right) \\
\end{align*}
and it is easy to check that:
\begin{align*}
    \cmmt{L_3}{X_q} &= \hbar q X_q \\
    \cmmt{L_{\pm}}{X_q} &= \hbar \sqrt{(1 \mp q)(1 \pm q + 1)} X_{q \pm 1}
\end{align*}
We see these are the same as the standard formulae for the actions of $L_3$ and $L_{\pm}$ on angular momentum states $\ket{1~q}$. This implies $X_q \ket{n~l~m}$ behaves exactly like the product states $\ket{1~q}\ket{l~m}$ with respect to angular momentum. This means:
\begin{equation*}
    \vec L \left(X_q \ket{n~l~m}\right) = \cmmt{\vec L}{X_q}\ket{n~l~m} + X_q \left(\vec L \ket{n~l~m}\right)
\end{equation*}
corresponds to:
\begin{equation*}
    \vec L \left(\ket{1~q}\ket{l~m}\right) = \left(\vec L \ket{1~q}\right)\ket{l~m} + \ket{1~q}\left(\vec L\ket{l~m}\right)
\end{equation*}
By the standard results for adding angular momentum, we can deduce that $X_q \ket{n~l~m}$ has angular momentum numbers:
\begin{align*}
    l+1, l, \text{ or } l-1&\text{ for } \vec L^2 \\
    m + q &\text{ for } L_3
\end{align*}
hence for $\opbraket{n'~l'~m'}{X_q}{n~l~m}$ to be non-zero we need $l' = l$ or $l \pm 1$, and $m' = m + q$. This is the second set of selection rules. Combining this with the first excludes $l = l'$ and so we have $l' = l \pm 1$.

In summary, the possible atomic transitions are given by:
\begin{align*}
    l' &= l \pm 1 \\
    m' = 
    \begin{cases}
        m & \text{ if } \EField_3 \neq 0 \\
        m \pm 1 & \text{ if } \EField_1 \neq 0 \text{ or } \EField_2 \neq 0
    \end{cases}
\end{align*}
and in addition we must have $E_{n'} = E_n \pm \hbar \omega_0$.
\end{document}